%%% Technical appendix for the AAMAS 2027 submission (anonymous). Supplementary material, not page-limited.
\documentclass[11pt]{article}
\usepackage[margin=1in]{geometry}
\usepackage[T1]{fontenc}
\usepackage{mathptmx}
\usepackage{amsmath,amssymb,amsthm}
\usepackage{booktabs,array}
\usepackage{algorithm,algorithmic}
\usepackage[hidelinks]{hyperref}
\usepackage{xr}
\externaldocument{main}
\usepackage{url}
\sloppy
\emergencystretch=3em
\urlstyle{same}
\renewcommand{\thesection}{\Alph{section}}
\newtheorem{theorem}{Theorem}[section]
\newtheorem{lemma}[theorem]{Lemma}
\newtheorem{proposition}[theorem]{Proposition}
\newtheorem{corollary}[theorem]{Corollary}
\theoremstyle{definition}
\newtheorem{definition}[theorem]{Definition}
\newcommand{\ind}{\mathbf 1}
\newcommand{\red}[1]{\textcolor{red}{#1}}
\usepackage{xcolor}
\title{Subjective SD-EF1 and Market EF1:\\Technical Appendix}
\author{Anonymous Authors}
\date{}
\begin{document}
\maketitle

\noindent This appendix supplements the AAMAS 2027 submission ``Subjective SD-EF1 and Market EF1: Two-Agent Existence and Barriers for Three Agents''. Theorem, lemma and section numbers refer to that paper (``main paper'' below); the appendix sections are lettered. Nothing here extends the finite existence theorem (Theorem~\ref{thm:twelve}) to arbitrary size. The code, logs and data are in the anonymous repository \url{https://anonymous.4open.science/r/price-and-prize-0625}; each result below names the directory that reproduces it.

\medskip\noindent\textbf{Conventions.} Agents in the three-agent setting are $0,1,2$; goods are numbered in market order, $\sigma_0=\pi=(0,\dots,m-1)$; an allocation is written $(D,S,X)=(A_0,A_1,A_2)$ and recorded as an owner vector $o$ with $o_g$ the recipient of $g$. Subjective SD-EF1 is the prefix condition $|A_j\cap P_i(t)|\le|A_i\cap P_i(t)|+1$ for all $i,j,t$. For $m=3q$, market-block balance gives every agent one good in each market triple. Padding adds universally bottom-ranked virtual goods; deleting them preserves all original prefixes, full blocks and short-block conditions (it is not asserted that every unpadded solution extends to a padded one).

\tableofcontents

%%%%%%%%%%%%%%%%%%%%%%%%%%%%%%%%%%%%%%%%%%%%%%%%%%%%%%%%%%%%%%%%%%%%%%%
\section{Two Agents: Chores, Mixed Items and Checks}\label{app:two}
\paragraph{Chores.}
Let two agents rank chores by decreasing subjective cost, and let $c$ be an identical inclusion-monotone market cost. \emph{Claim:} some allocation satisfies, for $i\ne j$, $|A_i\cap P|\le|A_j\cap P|+1$ for every prefix $P$ of agent~$i$'s ranking (subjective SD-EF1 for chores) and, whenever $A_i\ne\varnothing$, $c(A_i\setminus\{g\})\le c(A_j)$ for some $g\in A_i$ (market EF1 for chores); it is found with $O(\log m)$ value queries. \emph{Proof.} Rank the chores by decreasing cost, apply Theorem~\ref{thm:two} with $c$ as a goods valuation to obtain $(B_1,B_2)$, and return $(B_2,B_1)$. The goods inequality $|B_j\cap P|\le|B_i\cap P|+1$ becomes the chore inequality $|A_i\cap P|\le|A_j\cap P|+1$, and $c(B_i)\ge c(B_j\setminus\{g\})$ becomes $c(A_i\setminus\{g\})\le c(A_j)$. \qed

\paragraph{Mixed items with a common classification.}
Let $M=G\sqcup C$ be a classification into goods and chores shared by both agents and the market. Agent~$i$ has magnitudes $w_i\ge0$, with $u_i(x)=w_i(x)$ for $x\in G$ and $u_i(x)=-w_i(x)$ for $x\in C$. We assume that agent~$i$ reports one strict ranking of \emph{all} items by magnitude. Mixed EF1 for $i$ towards $j$ holds if $u_i(A_i)\ge u_i(A_j)$ or, after removing one good from $A_j$ or one chore from $A_i$, the inequality holds. An allocation is \emph{mixed SD-EF1} if it is mixed EF1 for every additive magnitude vector consistent with each agent's ranking. The market utility $f$ is arbitrary except that adding a good never decreases it and adding a chore never increases it; mixed EF1 for $f$ is defined analogously.

\begin{theorem}\label{app:mixed}
Some allocation is mixed SD-EF1 and mixed EF1 for $f$. It is found with $O(\log m)$ value queries.
\end{theorem}
\begin{proof}
Define the benefit bundles $B_i=(A_i\cap G)\cup(A_j\cap C)$; they partition $M$, the map $A\mapsto B$ is a bijection, and $u_i(A_i)-u_i(A_j)=w_i(B_i)-w_i(B_j)$. Run the reversal path of the main paper on the magnitude rankings in benefit coordinates.

\emph{Subjective side.} If $x$ is the $\sigma_i$-first item of $B_j$, the SD-EF1 prefix inequalities give $w_i(B_i)\ge w_i(B_j\setminus\{x\})$ for every consistent $w_i$. Here $x$ is either a good of $A_j$ or a chore of $A_i$, and removing it raises $u_i(A_i)-u_i(A_j)$ by exactly $w_i(x)$, so it is a mixed-EF1 repair.

\emph{Market side.} Put $F(T)=f(T\triangle C)$. Adding a good to $T$ adds a good to $T\triangle C$, and adding a chore to $T$ removes a chore from $T\triangle C$, so $F$ is inclusion-monotone, and $F(B_i)=f(A_i)$. For $x\in B_j$ the mixed-EF1 repairs available to agent~$i$ are, in benefit coordinates,
\begin{gather*}
 F(B_i)\ge F(B_j\setminus\{x\})\quad(x\in G),\qquad F(B_i\cup\{x\})\ge F(B_j)\quad(x\in C),
\end{gather*}
together with $F(B_i)\ge F(B_j)$: indeed $(B_j\setminus\{x\})\triangle C=A_j\setminus\{x\}$ for a good $x$ of $A_j$, and $(B_i\cup\{x\})\triangle C=A_i\setminus\{x\}$ for a chore $x$ of $A_i$. Let $U$ ($L$) contain the allocations in which agent~2 (agent~1) fails market mixed EF1. Failure forces $F(B_2)<F(B_1)$ in $U$ and $F(B_1)<F(B_2)$ in $L$, so $U\cap L=\varnothing$.

\emph{No jump.} Consider a benefit swap $(X+g,Y+h)\to(X+h,Y+g)$. If the first state is in $U$ and the second in $L$, use the failed repair with $g$, which lies in the envied bundle both times: for $g\in G$, $F(X)>F(Y+h)\ge F(Y)>F(X+h)\ge F(X)$; for $g\in C$, $F(X+g)>F(Y+h+g)\ge F(Y+g)>F(X+h+g)\ge F(X+g)$. If the first state is in $L$ and the second in $U$, use $h$: for $h\in G$, $F(Y)>F(X+g)\ge F(X)>F(Y+g)\ge F(Y)$; for $h\in C$, $F(Y+h)>F(X+g+h)\ge F(X+h)>F(Y+g+h)\ge F(Y+h)$. Each is a contradiction. For a move $(X+g,Y)\to(X,Y+g)$ from $U$ to $L$ the repair with $g$ gives $F(Y)<F(X)$ and $F(X)<F(Y)$ if $g\in G$, and $F(Y+g)<F(X+g)$ and $F(X+g)<F(Y+g)$ if $g\in C$; from $L$ to $U$ the unrepaired envy gives $F(X+g)<F(Y)$ and $F(Y+g)<F(X)$, which contradict monotonicity. The endpoints of the path are reversals, so exactly as in the main proof some state lies in neither $U$ nor $L$. Bisection on the sign of $F(B_1)-F(B_2)$ finds adjacent states with opposite signs; if the transferred item $g$ is a good, one query of $F(X)$ decides as in the main proof, and if it is a chore, one query of $F(Y+h+g)$ decides between the repair $F(Y+h+g)\ge F(X+g)$ at the earlier state and $F(X+h+g)\ge F(X+g)>F(Y+h+g)\ge F(Y+g)$ at the later one.
\end{proof}
A benefit swap of oppositely signed items transfers both original items in one direction; nothing is claimed about arbitrary swaps in original coordinates. Without a common classification the statement can fail: with subjective utilities $(2,-1)$ for both agents and $v(S)=|S|$, the only subjectively EF1 allocations put both items together, and neither is market EF1.

\paragraph{Checks.}
The tests of the two-agent construction (directories \path{verification/monotone_and_selection} and \path{verification/queries_and_defects}) cover: every second ranking for $m\le8$ with the first fixed (46{,}233 normalized profiles, 230{,}256 path allocations); all 168 monotone Boolean functions on four goods with all normalized ranking pairs (4{,}032 profile--valuation instances); $3{,}000$ random monotone-table instances each for goods, chores and mixed items with full allocation enumeration ($9{,}000$ in total, about $2\cdot10^6$ allocations enumerated); $3{,}000$ monotone-oracle instances with $11\le m\le60$; and $3{,}000$ instances forcing opposite failing endpoints. A second package tests the logarithmic-query algorithm on $2{,}000$ goods, $2{,}000$ chore and $2{,}000$ mixed instances with $m\le64$ (market families: additive, XOS, budget-additive, coverage, concave, supermodular, arbitrary monotone tables), with $49{,}984$ value queries in total, at most $13$ per instance and always within the proved bound, and $4{,}648$ explicit repair certificates. An independently written implementation (\path{src/verify_monotone_ours.py}) ran $3{,}000$ random ranking profiles with up to $24$ goods against $16{,}000$ random monotone valuations: none of $19{,}715$ path states violated subjective SD-EF1, every valuation found a market-EF1 state, and no step jumped between $U$ and $L$. Appendix~\ref{app:family} adds exact linear-programming checks.

%%%%%%%%%%%%%%%%%%%%%%%%%%%%%%%%%%%%%%%%%%%%%%%%%%%%%%%%%%%%%%%%%%%%%%%
\section{Cover Sizes for Two Agents}\label{app:family}
This appendix supports Section~\ref{sec:family} of the main paper. Code: \path{src/n2_cover.py}; results: \path{results/n2_cover_experiments.json}. All statements here are computations on sampled instances, except Proposition~\ref{prop:size} of the main paper, which is proved there.

\paragraph{Definitions.}
An allocation of $m$ goods between two agents is a $\pm1$ vector $x$ ($x_g=+1$ iff agent~1 holds $g$). Along a ranking $\sigma$, agent~1 is subjectively SD-EF1 iff $\min_t\sum_{s\le t}x_{\sigma(s)}\ge-1$, agent~2 iff $\max_t\sum_{s\le t}x_{\sigma(s)}\le1$. $S$ is the set of allocations satisfying both (for $\sigma_1,\sigma_2$ respectively). The market order is the identity for additive valuations. A profile $(\sigma_1,\sigma_2)$ is \emph{hard} if no allocation of $S$ has all market prefix differences $|\sum_{s\le t}x_s|\le1$, i.e.\ none is market-block-balanced, so no single allocation is EF1 for every additive valuation. Hardness is only a sampling filter for the monotone experiments. An allocation and its agent-swap behave identically for every common market valuation.

\paragraph{Deciding whether a family covers.}
For a family $\mathcal F$ and a class of valuations, $\mathcal F$ fails iff some valuation of the class makes every member violate market EF1. For each member, the violation is one of two types (agent~2 envies beyond one good, or agent~1 does), each a system of \emph{strict linear inequalities} in the valuation: for additive nonincreasing $v$ the type-$U$ inequality is $\sum_g(x_g-\mathbf 1[g=g_1])v_g>0$ with $g_1$ the first good of agent~1's bundle, and symmetrically for type~$L$; for arbitrary monotone $v$ (a vector indexed by the $2^m$ bundles with $v(\varnothing)=0$ and $v(T)\le v(T\cup g)$) the type-$U$ condition is $v(A_1\setminus g)-v(A_2)\ge1$ for all $g\in A_1$ after scaling, and symmetrically. $\mathcal F$ fails iff one of the $2^{|\mathcal F|}$ type patterns is feasible; each pattern is one linear program (HiGHS). A feasible pattern yields a witness valuation that is checked directly against the definition.
The linear programs were validated against direct evaluation of the EF1 definition on $40{,}000$ sampled additive valuations (experiment E0, $m=10$): of $3{,}000$ random single allocations, $2{,}900$ were declared failing and $100$ declared covering, and none of the latter was broken by a sampled valuation; of the $25{,}292$ covering pairs found on $120$ hard profiles none was broken by a sampled valuation, while each of $3{,}000$ pairs declared non-covering was broken by some sampled valuation. For the reversal path (below), whose covering property is a theorem, the exact test agrees.

\paragraph{E1: two allocations for additive valuations.}
Table~\ref{tab:e1} shows that on every sampled hard profile (600 random profiles with $8\le m\le13$) some pair of members of $S$ covers all additive valuations. A hill-climb over profiles that minimizes the number of pairs of members of $S$ with disjoint defect masks (a necessary condition for a covering pair) never produced a profile without a covering pair; the reported best profiles still have many candidate pairs.
% generated by gen_tables.py from results/n2_cover_experiments.json; do not edit by hand
\begin{table}[h]
\centering\small
\begin{tabular}{@{}rrrrrrr@{}}
\toprule
$m$ & hard profiles & random pairs drawn & without 2-cover & climb restarts & min.\ disjoint pairs & 2-cover there\\
\midrule
8 & 100 & 865{,}746 & 0 & 581 & 24 & yes\\
9 & 100 & 199{,}988 & 0 & 738 & 33 & yes\\
10 & 100 & 398{,}740 & 0 & 256 & 25 & yes\\
11 & 100 & 196{,}252 & 0 & 237 & 23 & yes\\
12 & 100 & 379{,}282 & 0 & 56 & 39 & yes\\
13 & 100 & 246{,}495 & 0 & 49 & 41 & yes\\
\bottomrule
\end{tabular}
\caption{Additive markets, two agents (E1). For each $m$, 100 random hard profiles were drawn and a covering pair of subjectively SD-EF1 allocations was searched among all pairs with disjoint defect masks (necessary) and verified by linear programming. The hill-climb minimizes the number of pairs with disjoint defect masks over profiles; the last two columns report the best profile it found (fewest such pairs) and whether it still has a covering pair. No profile without a covering pair was found.}
\label{tab:e1}
\end{table}


\paragraph{E2: exact minimum covers for monotone valuations.}
The minimum size of a cover of all monotone valuations consisting of members of $S$ (agent-swapped twins identified) is found by counterexample-guided search: the master problem is a minimum hitting set over witness valuations (initially $120$ random monotone valuations of five families, then witnesses returned by the linear programs), the check is the linear program above. It stops when the proposed family is a cover, hence optimal. On all $20$ sampled hard profiles with $m\in\{7,8\}$ the minimum is exactly $\lceil m/2\rceil=4$ (Table~\ref{tab:e2}). We did not obtain exact minima for $m\ge9$: the counterexample-guided search then needs many more witnesses and did not converge within our time budget, which is why E3 below only checks the path.
% generated by gen_tables.py from results/n2_cover_experiments.json; do not edit by hand
\begin{table}[h]
\centering\small
\begin{tabular}{@{}rrrll@{}}
\toprule
$m$ & profiles & $\lceil m/2\rceil$ & minimum cover size (count) & $|\mathcal S|$\\
\midrule
7 & 8 & 4 & 4: 8 & 16--22\\
8 & 12 & 4 & 4: 12 & 16--22\\
\bottomrule
\end{tabular}
\caption{Exact minimum cover of all inclusion-monotone valuations by subjectively SD-EF1 allocations (E2), on sampled hard profiles; $|\mathcal S|$ is the number of subjectively SD-EF1 allocations of the profile.}
\label{tab:e2}
\end{table}


\paragraph{E3: the reversal path is an exact cover, and minimal.}
For each sampled hard profile we built the path of Theorem~\ref{thm:two} with an independent implementation, checked that all its members lie in $S$, and verified with exact linear programs that (a) the path is a cover of all monotone valuations and (b) removing any of its members (up to swap twins) leaves a family that fails (Table~\ref{tab:e3}). Since the path's last allocation is the swap of its first, a path of $\lceil m/2\rceil+1$ allocations has $\lceil m/2\rceil$ market classes, matching the minimum found in E2 for $m\le8$. Minimality is weaker than minimum: it says the path cannot be shortened by deleting members, not that no other family is smaller.
% generated by gen_tables.py from results/n2_cover_experiments.json; do not edit by hand
\begin{table}[h]
\centering\small
\begin{tabular}{@{}rrrrrr@{}}
\toprule
$m$ & profiles & market classes on the path & path $\subseteq S$ & path is a cover & every member necessary\\
\midrule
9 & 4 & 5 & 4/4 & 4/4 & 4/4\\
10 & 4 & 5 & 4/4 & 4/4 & 4/4\\
11 & 4 & 6 & 4/4 & 4/4 & 4/4\\
\bottomrule
\end{tabular}
\caption{The reversal path of Theorem 3.1 on sampled hard profiles (E3), verified by exact linear programs over all monotone valuations. ``Market classes'' counts the path's allocations up to exchanging the two bundles.}
\label{tab:e3}
\end{table}


\paragraph{E4: the disjoint-defect route fails for two agents.}
For two agents call an allocation \emph{near-block-balanced} if all market prefix differences are at most two in absolute value; then defects are the prefixes with difference $\pm2$. For the eleven-good profile $\sigma_1=(0,3,4,7,2,6,1,9,8,10,5)$, $\sigma_2=(7,3,4,6,0,2,8,9,10,5,1)$ (identity market order) there are $81$ subjectively SD-EF1 allocations, none market-block-balanced, $24$ near-block-balanced, and no two of the $24$ have disjoint defect sets; yet some pair of allocations (one not near-block-balanced) covers all additive valuations. Exhaustive enumeration ($2^{11}$ allocations) and linear programs give these counts (function \texttt{experiment\_E4} of \path{src/n2_cover.py}).

\paragraph{E5: covering pairs off the reversal path.}
Table~\ref{tab:e5} shows that a covering pair of additive valuations is not always found among the reversal path's allocations: for $60$ sampled hard profiles per $m$, the path contains a covering pair in $57$, $47$, $49$ and $30$ cases for $m=8,9,10,11$. For the other profiles a pair exists among other members of $S$ (E1). A proof that two allocations always suffice for additive valuations would therefore need a construction different from the path.
% generated by gen_tables.py from results/n2_cover_experiments.json; do not edit by hand
\begin{table}[h]
\centering\small
\begin{tabular}{@{}rrr@{}}
\toprule
$m$ & sampled hard profiles & with a covering pair among the path's allocations\\
\midrule
8 & 60 & 57\\
9 & 60 & 47\\
10 & 60 & 49\\
11 & 60 & 30\\
\bottomrule
\end{tabular}
\caption{Additive markets (E5): a covering pair of subjectively SD-EF1 allocations need not lie on the reversal path.}
\label{tab:e5}
\end{table}


\paragraph{Reproduction.}
\texttt{python n2\_cover.py all} in \path{src/} (about eight minutes on 12 cores) recomputes E1--E5 with fixed seeds and writes \path{results/n2_cover_experiments.json}; \path{paper/gen_tables.py} turns that file into the tables above.

%%%%%%%%%%%%%%%%%%%%%%%%%%%%%%%%%%%%%%%%%%%%%%%%%%%%%%%%%%%%%%%%%%%%%%%
\section{Coverage Computations for Three Agents}\label{app:coverage}
\subsection{Twelve goods}
Relabel goods by their position in $\sigma_1$. Consider every partition $Q$ of $\{0,\dots,11\}$ into four unordered triples; $\sigma_1$ is the identity. The market order itself is not the identity after this normalization: $Q$ retains exactly the market information needed for block balance, since agent~0 accepts every rainbow allocation of $Q$ independently of the order of its blocks and the order inside each block.

\paragraph{Canonical partitions and candidates.}
Choose the smallest remaining good $a$, enumerate its two partners, remove that triple, and recurse; each partition is generated once, giving $12!/((3!)^44!)=15{,}400$. For each $Q$ enumerate all six owner permutations per triple and retain the candidates that give agent~1 one good of every identity-order triple and satisfy all of its prefix conditions; this is $\mathcal F(Q)$. The code checks fairness separately from own-bundle balance: taking one good from each own-ranking triple does not by itself protect the owner against both opponents when the opponents are not block-balanced.

\paragraph{Coverage recursion.}
Let $T$ be the set of exposed goods in a prefix of agent~2's ranking and $L$ the candidates that satisfy all exposed prefix conditions in their encountered order. Algorithm~\ref{alg:cover} decides whether every continuation is covered; the empty-family test precedes any acceptance rule, and at a fully exposed prefix any surviving allocation contains four agent-2 goods and triggers acceptance.

\begin{algorithm}[h]
\caption{\textsc{Cover}$(T,L)$ for a fixed partition $Q$}\label{alg:cover}
\begin{algorithmic}[1]
\IF{$L=\varnothing$} \RETURN false \ENDIF
\IF{some $A\in L$ has $|A_2\cap T|\ge3$} \RETURN true \ENDIF
\IF{$(T,L)$ is a memoized covered state} \RETURN true \ENDIF
\FOR{each $g\in M\setminus T$}
  \STATE $T'\gets T\cup\{g\}$; $L'\gets\{A\in L: |A_j\cap T'|\le |A_2\cap T'|+1\text{ for }j=0,1\}$
  \IF{\textsc{Cover}$(T',L')$ is false} \RETURN false \ENDIF
\ENDFOR
\STATE memoize $(T,L)$ as covered; \RETURN true
\end{algorithmic}
\end{algorithm}

\begin{proposition}
For every reachable state $(T,L)$, Algorithm~\ref{alg:cover} returns true iff every ordering of $M\setminus T$ completes the exposed prefix to a ranking accepted by at least one candidate of $L$.
\end{proposition}
\begin{proof}
An empty $L$ admits no acceptable continuation. If a surviving candidate has three agent-2 goods, its count never decreases while every opposing bundle has total size four, so all future differences are at most one; this is the safe-acceptance case. For every other state, the continuations are partitioned by their next good $g$; at the new prefix the surviving candidates are exactly $L'$. Induction on $|M\setminus T|$ shows that all continuations are covered exactly when all recursive children are covered. No later feasibility depends on the order before $T$ except through survival in $L$, so memoizing the complete pair $(T,L)$ preserves the answer.
\end{proof}
The C++ implementation represents candidates by three good masks and $L$ by an integer bitset, precomputes validity and safe-acceptance masks for every prefix set, and keys its memo table by the prefix mask and the entire candidate bitset (equality does not rely on a hash fingerprint). Children with an immediate safe witness may be skipped; the order of the other children changes the traversal, not coverage.
\begin{center}\small
\begin{tabular}{@{}lr@{}}
\toprule
Quantity & Complete C++ run\\ \midrule
Canonical partitions & 15{,}400\\
Uncovered partitions & 0\\
Recursive calls & 70{,}831{,}414\\
Memoized covered states & 43{,}515{,}662\\
Memo hits & 27{,}315{,}752\\
Safe-prune events & 260{,}293{,}476\\
Max.\ calls per partition & 4{,}989\\
Candidate-family range & 64 / 1{,}296\\
\bottomrule
\end{tabular}
\end{center}
(Directory \path{verification/coverage12}; safe-prune events include children accepted without a recursive call.)

\paragraph{Independent checks and negative control.}
A second C++ checker, written separately (\path{src/cover_ours_general.cpp}, compiled with \texttt{-DMGOODS=12} and run with \texttt{OWN=1}), enumerates the same partitions and candidates but shares no code with the first and uses no safe-prune ordering; it covers all $15{,}400$ partitions and reproduces exactly the candidate-family range ($64$ to $1{,}296$) and the number of memoized covered states ($43{,}515{,}662$), while its number of recursive calls differs ($331{,}124{,}890$), as expected from the different traversal (about a minute on one core; \path{results/cover12_ours/}). A Python implementation builds explicit owners and checks prefixes separately; forty complete coverage problems agree with C++ on coverage and candidate counts. For a negative control require all three bundles to be rainbow on every identity-ranking triple, rather than only agent~1's; for
$Q=\{\{0,2,9\},\{1,7,8\},\{3,10,11\},\{4,5,6\}\}$ the stronger family has $48$ allocations and the routine returns the uncovered prefix $(0,2,4,1,3,5,6,9)$; the independent C++ checker also reports this partition (number $4709$ in the enumeration) as uncovered, and $12{,}150$ of the $15{,}400$ partitions under this control. This demonstrates rejection when the family is overconstrained; it is not a counterexample to the theorem.

\subsection{Fifteen goods}
Pad to $M=3q$ goods with $q=5$; let $T_0,\dots,T_{q-1}$ be the consecutive triples of $\sigma_1$, put $c_j=|S\cap T_j|$, $h_k=\sum_{j<k}(c_j-1)$ and $b(S)=\sum_j\max(c_j-1,0)$; own-bundle balance is $b(S)=0$. By Appendix~\ref{app:fifteen} every market-block-balanced subjectively SD-EF1 allocation has $h_k\ge0$ and $h_q=0$, so $c$ decomposes into $b(S)$ unit shifts, each moving one good of agent~1 from a later triple to an earlier one.

The recursion of the previous subsection is extended in two exact ways. First, a \emph{covering core} $K\subseteq L$ at $T$ is a subfamily that already covers every continuation; the union of the children's cores is a core of the parent, and a stored core $K$ may be reused at any later state $(T,L')$ with $K\subseteq L'$. Second, with $\sigma_1$ normalized to the identity the candidate families are invariant under permuting the first two labels and the final triple and, when the first $\sigma_1$ triple is additionally required to give one good to each agent, under permuting the first and the final triple; restricted runs that fail fall back to the smaller group. Over all $15!/((3!)^55!)=1{,}401{,}400$ normalized partitions ($82{,}670$ orbits of the larger group, $15{,}937{,}540{,}828$ recursive calls, $101{,}840$ coverage problems) every partition is covered by own-bundle-balanced or one-shift candidates. A rerun on a second machine reproduces every deterministic counter and all $82{,}670$ per-orbit records (directory \path{verification/coverage15}); an accounting script reconstructs every orbit and its weight (orbit sizes sum to $1{,}401{,}400$) and audits the integer logs, but is not a replay of the coverage states.

%%%%%%%%%%%%%%%%%%%%%%%%%%%%%%%%%%%%%%%%%%%%%%%%%%%%%%%%%%%%%%%%%%%%%%%
\section{Own-Bundle Balance and Thirteen Goods}\label{app:fifteen}
\begin{lemma}
Every market-block-balanced subjectively SD-EF1 allocation has $h_k\ge0$ for all $k$ and $h_q=0$.
\end{lemma}
\begin{proof}
At the first $3k$ goods of $\sigma_1$ let agent~1 hold $s$ goods; each opponent holds at most $s+1$, so $3k\le3s+2$ and $s\ge k$. All bundles have size $q$.
\end{proof}

\paragraph{The thirteen-good profile.}
Take $\sigma_0=(0,\dots,12)$, $\sigma_1=(0,1,3,12,2,6,9,10,4,7,5,8,11)$ and $\sigma_2=(3,2,0,1,6,5,12,4,8,11,7,9,10)$. Relabel the goods by their position in $\sigma_1$; the market blocks become $Q_A=\{0,1,4\}$, $Q_B=\{2,8,10\}$, $Q_C=\{5,9,11\}$, $Q_D=\{6,7,12\}$, $Q_E=\{3\}$, and $\sigma_2$ becomes $(2,4,0,1,5,10,3,8,11,12,9,6,7)$.

\begin{proposition}
No market-block-balanced subjectively SD-EF1 allocation of this profile is own-bundle balanced.
\end{proposition}
\begin{proof}
Own-bundle balance gives $|S|=4+\ind[12\in S]$ and market balance gives $|S|=4+\ind[3\in S]$, so $3\in S\iff12\in S$. The two block systems then leave three forms: (I) $S=\{a,3,8,v,12\}$ with $a\in\{0,1\}$, $v\in\{9,11\}$; (II) $S=\{2,4,u,v\}$ with $u\in\{6,7\}$, $v\in\{9,11\}$; (III) $S=\{a,5,u,10\}$ with $a\in\{0,1\}$, $u\in\{6,7\}$. (If $3,12\in S$ then $Q_D$ excludes $6,7$ and forces $8\in S$; otherwise $S$ takes $6$ or $7$, excluding $8$, and its $Q_B$-good is $2$ or $10$.) In case I, the first eleven goods of $\sigma_2$ contain all five goods of $S$ but at most three of $X$, because agent~2's $Q_D$-good is $6$ or $7$. In case II, the first two goods of $\sigma_2$ both belong to $S$. In case III, the first six goods of $\sigma_2$ contain three $S$-goods and agent~2's $Q_A$-good, so $2\in X$. The first five goods of $\sigma_1$ contain one $S$-good and two $X$-goods, so $3\notin X$, hence $3\in D$ and $|X|=4$. The first eleven goods of $\sigma_2$ contain three $S$-goods, so all four $X$-goods lie there and $12\in X$. The first nine goods of $\sigma_1$ contain four $D$-goods and three $S$-goods, so $9\notin D$, i.e.\ $9\in X$. The first eight goods of $\sigma_2$ have counts $(|D\cap P|,|S\cap P|,|X\cap P|)=(3,3,2)$, so the next good $11$ must also belong to $X$; then $X$ contains both $9$ and $11$ from $Q_C$, contradicting market balance.
\end{proof}
\begin{center}\small
\begin{tabular}{@{}rrrr@{}}
\toprule
$m$ & Balanced & Fair & Fair, own-balanced\\ \midrule
13 & 3{,}888 & 206 & 0\\
15 (padded) & 7{,}776 & 412 & 0\\
\bottomrule
\end{tabular}
\end{center}
Enumeration of the thirteen-good profile and of its padding by two universally bottom-ranked goods; every fair padded allocation has $b(S)=1$. The counts were reproduced by an independent implementation (\path{src/verify_thirteen.py}). The example allocation $A_0=\{0,3,7,9\}$, $A_1=\{1,4,6,10\}$, $A_2=\{2,5,8,11,12\}$ of the main paper is fair.

\paragraph{No constant number of shifts.}
Concatenate $h$ copies of the padded profile, contiguous in the market and in both real rankings. By the module factorization of Appendix~\ref{app:distance}, every fair allocation restricts to fair allocations of the copies, and triples of $\sigma_1$ do not cross copies, so $b(S)\ge h$; concatenating one witness per copy attains $h$. The maximum height $\max_kh_k$ of these witnesses is $1$, so a bound on the height, rather than on the number of shifts, is not excluded.

%%%%%%%%%%%%%%%%%%%%%%%%%%%%%%%%%%%%%%%%%%%%%%%%%%%%%%%%%%%%%%%%%%%%%%%
\section{Market-Aligned Existence for All Sizes}\label{app:allsizes}
\paragraph{Converse for the one-singleton rule.}
Let $m=3q$ and suppose a prefix $P$ of $\sigma_2$ meets at least two triples in exactly one good. Choose one good of $P$ in every triple meeting $P$ and put these goods first in $\sigma_1$. Agent~1 takes the chosen good of every such triple. A triple with two goods in $P$ contains agent~2's two favourites of the triple; one is taken by agent~1 and the other by agent~2, so both count in $P$. A triple with three goods in $P$ contributes one to each. A singleton triple contributes one to agent~1 and none to agent~2. Hence $|A_1\cap P|-|A_2\cap P|\ge2$ and agent~2 is not SD-EF1; this certifies failure of the rule, not nonexistence.

\paragraph{Covered perturbations.}
If $\sigma_2$ arises from $\pi$ by disjoint adjacent transpositions, every prefix ends inside at most one triple that it meets in exactly one good. More generally, permute the triples and the goods within them and then exchange the two goods at selected triple boundaries: whenever the next triple starts early, the preceding triple already contributes two goods.

\paragraph{One arbitrary ranking with block-contiguous others.}
Call a ranking \emph{block-contiguous} if the goods of each market block are consecutive in it, in any block order and any internal order.
\begin{theorem}
Let $n\ge2$ and $m=qn+r$ with $0\le r<n$. If every agent except possibly one agent $f$ has a block-contiguous ranking, a market-block-balanced subjectively SD-EF1 allocation exists and is computed in $O(nm)$ time.
\end{theorem}
\begin{proof}
Let $T$ be the short block. Give its goods bijectively to a set $R$ of $r$ agents other than $f$, and let $H$ be the remaining agents other than $f$. In every full block, agent $f$ picks its favourite good, then the agents of $H$, then those of $R$, each taking its favourite remaining good. Agent $f$'s good in each full block precedes every other good of that block in $\sigma_f$, so $f$ dominates every opponent's full-block goods, and each opponent has at most one good of $T$. For an agent with a block-contiguous ranking, completed full blocks contribute equally to all agents, a completed short block contributes at most one good to anyone and one to the agent if it lies in $R$, and the current block contributes at most one good to each agent. An agent of $H$ picks before every agent of $R$ in each full block, which covers comparisons with $R$. In every case the lead of an opponent is at most one.
\end{proof}
For three agents with $\sigma_0=\pi$ and $3\mid m$, a block-contiguous ranking is one-singleton, so this case is contained in Theorem~\ref{thm:onesingleton} after exchanging agents~1 and~2 if needed; the theorem above adds short final blocks and $n>3$.

\paragraph{Modules.}
To find common modules, use a dynamic program on market boundaries that tests widths $3,6,9,12$. An interval $[a,b)$ is contiguous in ranking $i$ exactly when $\max_{a\le g<b}\operatorname{rank}_i(g)-\min_{a\le g<b}\operatorname{rank}_i(g)=b-a-1$. The widths are bounded, so these tests take constant time per option. In a successful segmentation each local solution assigns equal-sized bundles within a module. A subjective prefix contains complete modules, at most one partial module, and none of the others; all complete-module deficits cancel. This is why different module orders are allowed, and why arbitrary interleaving of their goods is not. For example, take market blocks $\{0,1,2\}$ and $\{3,4,5\}$ with owners $(0,1,2,0,1,2)$, $\sigma_1=(0,3,1,4,2,5)$ and $\sigma_2=(2,5,1,4,0,3)$. Each one-block restriction is fair, yet agent~1 receives nothing in its first two positions while agent~0 receives both goods: equal endpoint counts do not justify gluing chunks that interleave inside subjective prefixes.
For $m\le12$, choosing one of three goods per market block and then one of two residual goods gives at most $3^42^4=1{,}296$ candidates before rejection; for general $m$ this is only an exact search of the padded restricted family, and failure would not refute unrestricted market-aligned existence.

%%%%%%%%%%%%%%%%%%%%%%%%%%%%%%%%%%%%%%%%%%%%%%%%%%%%%%%%%%%%%%%%%%%%%%%
\section{Complexity Details}\label{app:complexity}
\paragraph{Cycle description.}
On $R=M\setminus S$ let $H$ combine the residual market pairs and the residual $\sigma_1$ pairs, with sides $(U_c,V_c)$ for each component $c$. For a $\sigma_2$ prefix $P$ set $r=|R\cap P|$, $s=|S\cap P|$ and $\delta_c(P)=|U_c\cap P|-|V_c\cap P|$. With sign $\epsilon_c=+1$ choosing $U_c$ for $X$, agent~2's remaining inequalities are
\[
 \sum_c\epsilon_c\delta_c(P)\ \ge\ \max\{2s-r-2,\,-(r\bmod2)\},
\]
which follows from $2x-r=\sum_c\epsilon_c\delta_c(P)$ and inequality~\eqref{eq:two} of the main paper ($x\ge\max\{s-1,\lfloor(t-s)/2\rfloor\}$ with $t-s=r$). Enumeration over the component signs takes $O(m2^c)$ time and $O(m)$ working space for $c$ components; this is an exact reduction, not an integrality theorem.

\paragraph{The reduction walk.}
With a positive anchor each clause gadget begins at $W=2$. If it has $t$ true literals the heights are:
\begin{center}\small
\begin{tabular}{@{}lrrrr@{}}
\toprule
True literals $t$ & 0 & 1 & 2 & 3\\ \midrule
After three literal goods & $-1$ & 1 & 3 & 5\\
After $\beta_j$ & $-2$ & 0 & 2 & 4\\
After complementary goods & 1 & 1 & 1 & 1\\
After $\alpha_{j+2}$ & 2 & 2 & 2 & 2\\
\bottomrule
\end{tabular}
\end{center}
Within the complementary segment there are $t$ negative steps; its height is at least $t-2$ for $t\ge1$. For $t=0$ all three steps are positive, starting from $-2$. The first bad prefix, when any clause is false, is therefore its $\beta_j$. In the anchor construction $\beta_j$ and $z_j$ share both blocks, so the suffix exchange of Lemma~\ref{lem:exchange} repairs all later bad prefixes at once. It does not change any variable-cycle truth choice and does not find a satisfying assignment for an unsatisfiable formula.

\paragraph{Linear relaxation.}
Consider the linear program with $w_g\in[0,1]$ for residual goods, $w_g=0$ on $S$, residual market-pair sums equal to one, and the prefix bounds~\eqref{eq:one}--\eqref{eq:two} with $x=\sum_{g\in P}w_g$. For a reset matching with the suffix restriction, giving every residual good weight $1/2$ is feasible: every pair sum is one; agent~1's reset-prefix constraints hold; before the suffix $x=r/2\ge\lfloor r/2\rfloor$, and during it $x=q\ge s-1$. Because the rounded lower bounds are themselves derived from integrality, this only shows that hardness is invisible to this particular relaxation. A twelve-good example has a genuine integrality gap: with
\begin{align*}
 \sigma_1&=(10,0,2,\ 1,3,5,\ 4,6,8,\ 7,9,11), &
 \sigma_2&=(10,1,0,3,6,9,4,7,\ 2,5,8,11),
\end{align*}
and $S=\{2,5,8,11\}$ the residual graph is one cycle with sides $U=\{10,1,4,7\}$ and $V=\{0,3,6,9\}$, and its $\sigma_2$ word is $UUVVVVUU$. Choosing $V$ fails at the first two goods, choosing $U$ at the first six. Fractionally the pair equalities force weight $p$ on $U$ and $1-p$ on $V$, and the two prefixes force $p\ge1/2$ and $p\le1/2$, so the unique relaxed completion is the half solution; yet a different matching works: $S'=\{2,5,8,9\}$, $X=U$, $D'=\{0,3,6,11\}$. Fixed-$S$ infeasibility is therefore not a counterexample to existence. With six goods, $\sigma_1=(0,1,2,3,4,5)$, $S=\{2,5\}$, $\sigma_2=(2,5,0,1,3,4)$, the first two positions have $s=2$ and $x=0$ for every fixed-$S$ fractional allocation, violating $x\ge s-1$: this instance is fractionally infeasible.

\paragraph{Linear-time applicability test of the exchange lemma.}
Compute all $\sigma_2$ prefix counts, the slack $h(t)=x(t)-s(t)$, and the first and last prefixes $f,\ell$ with $d(t)-x(t)=2$. If there are none, the input is already fair under the other premises. For a candidate pair with $b\in D$ and $a\in S$ the affected interval of prefix lengths is $[\operatorname{rank}_2(b)+1,\operatorname{rank}_2(a)]$ (zero-based positions). Reject it unless it contains $f$ and $\ell$ and unless $b$ precedes $a$ in the same $\sigma_1$ block. Every remaining interval contains $f$; precompute minima of $h$ extending left and right from $f$; the minimum on any such interval is the smaller of its two pivoted minima, so the slack hypothesis is tested in constant time and the whole test takes $O(m)$. Failure of this test only means that no exchange meeting these sufficient premises was found.

\paragraph{Integral negative controls.}
The contradictory formula with clauses $(x,x,x)$ and $(\neg x,\neg x,\neg x)$ produces a $30$-good profile with no fixed-matching completion, and the formula with all eight sign patterns on three variables a $102$-good profile with the same property; cycle enumeration rejects every fixed-matching choice, while both profiles admit the proved repair.

\paragraph{Exhaustive audit of the reduction.}
An exhaustive C++ audit enumerates every ordered three-literal formula with at most three clauses up to variable renaming ($10{,}840{,}297$ formulas: $1$, $40$, $12{,}992$ and $10{,}827{,}264$ with $0,1,2,3$ clauses; all signs, repeated occurrences, tautologies, literal orders and clause orders included), checks every raw residual-pair orientation ($346{,}680{,}644$ tests) against the original prefix conditions of agent~1, tests every survivor on $\sigma_2$ and compares each outcome with the truth table; there are no disagreements. An independent Python reference audit checks all $13{,}033$ formulas with at most two clauses and $1{,}000$ further three-clause formulas with the same result. (Directory \path{verification/repair}; we re-ran the C++ audit and it reproduces the shipped counts.)

%%%%%%%%%%%%%%%%%%%%%%%%%%%%%%%%%%%%%%%%%%%%%%%%%%%%%%%%%%%%%%%%%%%%%%%
\section{The Repair-Distance Family}\label{app:distance}
Let $q=2r$, index market triples by $\mathbb Z_q$, and write $a_i=3i$, $b_i=3i+1$, $z_i=3i+2$, so that market triple $i$ is $(a_i,b_i,z_i)$. Set
\begin{align*}
\sigma_1&=\mathop{\mathrm{concat}}_{j=0}^{q-1}(a_{j-1},b_{j-2},z_j),\\
\sigma_2&=(a_0,a_1,b_0,b_1,b_2,b_3,\ a_2,\dots,a_{q-1},\ b_4,\dots,b_{q-1},\ z_0,\dots,z_{q-1}),
\end{align*}
and $S=\{z_0,\dots,z_{q-1}\}$, the bottom $q$ goods of $\sigma_2$, a reset matching. Market triple $i$ meets $\sigma_1$ triples $i$ ($z_i$), $i+1$ ($a_i$) and $i+2$ ($b_i$); these are distinct for $q\ge4$, and the $a$- and $z$-edges already connect all indices.

\paragraph{No completion.}
The residual market pairs are $(a_i,b_i)$ and the residual $\sigma_1$ pairs are $(a_i,b_{i-1})$, forming one alternating cycle with sides $\{a_i\}$ and $\{b_i\}$. By the reset lemma these are agent~2's only options. Giving it all $b_i$ fails at the first two goods of $\sigma_2$; giving it all $a_i$ fails at the first six, where agent~2 has two goods and agent~0 four.

\paragraph{Lower bound.}
Any other perfect matching induces a permutation $i\mapsto i+\delta_i$ of $\mathbb Z_q$ with $\delta_i\in\{0,1,2\}$. The increments of a nontrivial cycle are positive and sum to a multiple of $q$, so the cycle has at least $q/2=r$ elements, each replacing one member of $S$.

\paragraph{Attaining allocations.}
With own-bundle balance, $X=\{a_i\}$, $S_*=\{b_i:i\text{ even}\}\cup\{z_i:i\text{ odd}\}$ and $D=M\setminus(S_*\cup X)$ is fair with $|S\setminus S_*|=r$. Without it, $\widehat S=(S\setminus\{z_{q-1}\})\cup\{a_{q-1}\}$, $\widehat X=(\{a_i\}\setminus\{a_2,a_{q-1}\})\cup\{b_2,b_{q-1}\}$ and $\widehat D=M\setminus(\widehat S\cup\widehat X)$ is fair with one replacement. For $r=2$ the instance is $\sigma_1=(9,7,2,0,10,5,3,1,8,6,4,11)$, $\sigma_2=(0,3,1,4,7,10,6,9,2,5,8,11)$, $S=\{2,5,8,11\}$, with repair $\widehat S=\{2,5,8,9\}$, $\widehat X=\{0,3,7,10\}$, $\widehat D=\{1,4,6,11\}$.
\begin{center}\small
\begin{tabular}{@{}rrrrr@{}}
\toprule
$r$ & Balanced & Fair & Own-balanced & Min.\ replacements (own-bal.\ / any)\\ \midrule
2 & 1{,}296 & 211 & 27 & 2 / 1\\
3 & 46{,}656 & 4{,}571 & 207 & 3 / 1\\
4 & 1{,}679{,}616 & 96{,}276 & 1{,}341 & 4 / 1\\
\bottomrule
\end{tabular}
\end{center}
Exhaustive checks: market-block-balanced allocations, fair ones, fair ones with own-bundle balance, and minimum replacements; the $r=4$ row was computed independently of the supplied package (\path{src/verify_repair_onesingleton.py}), and the explicit repaired allocation is fair for $r=2,\dots,10$.

\paragraph{Modules.}
If modules consist of full market blocks and are contiguous in every real ranking, a market-block-balanced allocation is globally SD-EF1 iff it is SD-EF1 within every module: completed modules contribute equally, and every local prefix embeds into a global one. Replacement distance therefore adds over copies of the $r=2$ instance.

%%%%%%%%%%%%%%%%%%%%%%%%%%%%%%%%%%%%%%%%%%%%%%%%%%%%%%%%%%%%%%%%%%%%%%%
\section{Disjoint Defects and Selection Obstructions}\label{app:defect}
\paragraph{Proposition~\ref{prop:defect} of the main paper, and the ten-good profile.}
On the ten-good profile $\sigma_0=(7,4,2,5,0,8,9,6,3,1)$, $\sigma_1=(4,2,0,3,7,5,9,1,8,6)$, $\sigma_2=(4,0,3,2,7,5,9,1,8,6)$ with market order $0,\dots,9$, the $483$ subjectively fair allocations include $34$ near-block-balanced ones and none that is market-block-balanced. Their defect sets are $\{6\}$ (9 allocations), $\{3,6\}$ (3), $\{3,9\}$ (2), $\{6,9\}$ (12) and $\{3,6,9\}$ (8); the only disjoint type is $\{6\}$ with $\{3,9\}$ ($18$ pairs). Choosing the lexicographically earliest defect set (which has no partner), maximizing total subjective slack, or requiring both members to minimize the number of defects all fail, although a disjoint pair exists. Prescribing the parity class of defective boundaries also fails on a separate ten-good profile.

\paragraph{Twelve goods: no disjoint pair.}
For
\begin{align*}
\sigma_0&=(0,7,2,5,1,11,8,3,6,9,4,10),\quad
\sigma_1=(0,11,2,7,1,5,10,9,3,8,6,4),\\
\sigma_2&=(5,0,2,11,7,1,4,3,9,8,6,10)
\end{align*}
with market order $0,\dots,11$, enumeration of all $531{,}441$ allocations gives $1{,}080$ subjectively fair and $58$ near-block-balanced allocations, with defect sets $\{3,6\}$ (12), $\{3,9\}$ (6), $\{6,9\}$ (36), $\{3,6,9\}$ (4), no two of which are disjoint. The profile still has a two-allocation cover, $A=(1,0,2,1,2,2,0,0,2,1,0,1)$ and $B=(1,1,0,2,0,2,1,2,0,2,1,0)$ as owner vectors, whose positive-row envelopes sum to a nonpositive vector; $B$ is not near-block-balanced. This refutes the disjoint-defect route, not Open Question~1. (Independent enumeration: \path{src/verify_defect_counterexample.py}.)

\paragraph{Obstructions to keeping an initial partition.}
The following examples have $\sigma_0=(0,\dots,9)$ and initial classes $C_0=\{0,3,6,9\}$, $C_1=\{1,4,7\}$, $C_2=\{2,5,8\}$, each balanced in the market and in the displayed $\sigma_1$.
\emph{One within-block swap is insufficient:} for $\sigma_1=(0,5,4,9,1,8,2,3,7,6)$ and $\sigma_2=(5,4,8,9,3,6,0,7,1,2)$ all $60$ distinct allocations obtained by permuting the classes among the agents and applying at most one within-market-block transposition fail, although $57$ market-block-balanced fair allocations exist, $18$ of them own-bundle balanced, e.g.\ $D=\{0,3,7\}$, $S=\{2,4,6,9\}$, $X=\{1,5,8\}$.
\emph{Freezing an initial class is insufficient:} for $\sigma_1=(9,4,2,1,6,5,7,3,8,0)$ and $\sigma_2=(1,4,5,3,6,0,9,2,7,8)$, fixing any $C_i$ as the exact bundle $S$ and freely assigning the other goods to agents~0 and~2 subject to market-block balance leaves all $40$ candidates failing, while $78$ fair market-block allocations exist, eight own-balanced, including $D=\{0,5,8,9\}$, $S=\{2,3,6\}$, $X=\{1,4,7\}$.
\emph{Ten goods are minimal for these two rules:} for $m\le9$ every class of the supplied balanced partition has at most three goods; if every class has at most one good every assignment is fair for agent~2, otherwise choose the class whose second good appears first in $\sigma_2$: before that position each rival has at most one good, afterwards the chosen class has at least two and each rival at most three, so the class is acceptable to agent~2; agent~1 accepts either remaining class because the partition is balanced along $\sigma_1$. Minimality is for these rules with a supplied doubly balanced partition, not for all repair methods.
\emph{Full own-ranking balance can be too strong:} for $\sigma_1=(6,0,7,9,4,5,3,1,2,8)$ and $\sigma_2=(7,4,5,6,8,9,0,3,2,1)$ there are $87$ market-block-balanced fair allocations, of which $14$ retain own-bundle balance, but none keeps all three bundles balanced in $\sigma_1$. Parallel-edge exchanges preserve the multiset of owners in each $\sigma_1$ triple, so a sequence of only these exchanges starting from full balance cannot solve this instance; this limits that move set, not the market-aligned conjecture.

%%%%%%%%%%%%%%%%%%%%%%%%%%%%%%%%%%%%%%%%%%%%%%%%%%%%%%%%%%%%%%%%%%%%%%%
\section{Test Suites and Reproduction}\label{app:tests}
\paragraph{Finite existence.}
Additional tests cover $6{,}500$ instances with $m\le12$, $30{,}000$ restricted searches through $m=30$, and $6{,}200$ common-module instances. All outputs pass direct checks; successes above twelve goods do not extend Theorem~\ref{thm:twelve}.

\paragraph{Completion and repair.}
A separate checker evaluates the original market-block and subjective-prefix conditions, not the cycle reduction or the SAT semantics. The suites cover $11{,}300$ generated inputs (not claimed distinct profiles):
\begin{center}\small
\begin{tabular}{@{}p{.4\textwidth}rr@{}}
\toprule
Suite & Inputs & Internal cases\\ \midrule
Reset-cycle equivalence & 3{,}000 & 613{,}800\\
SAT reduction & 3{,}000 & 92{,}304\\
Small SAT, raw pairs & 300 & 109{,}200\\
Conditional exchange & 5{,}000 & ---\\
\bottomrule
\end{tabular}
\end{center}
Internal cases count raw pair orientations, truth/anchor assignments and raw pair orientations respectively; all checks pass. The reset suite finds $2{,}468$ feasible and $532$ infeasible prescribed matchings, all classified correctly by the cycle solver. For the SAT suite all $46{,}152$ positive-anchor assignments are repaired, including $25{,}263$ nontrivial exchanges; the rankings-only construction is checked on the same $3{,}000$ profiles. The general exchange suite returns $2{,}071$ fair allocations, including $41$ nontrivial repairs, and its applicability decision agrees with direct inspection of the lemma's hypotheses; it is not a universal repair success rate. Exact rational arithmetic checks the half solution. A $1{,}000$-clause, $12{,}006$-good reduction profile is constructed and directly checked, not exhaustively enumerated. The one-singleton rule passes $4{,}000$ random one-singleton profiles, and its converse construction defeats the rule on $3{,}791$ random profiles violating the condition.

\paragraph{Three-agent evidence.}
Notes and logs: \path{notes/n3_counterexample_search.md}, \path{results/}. Exhaustive enumeration over all $\binom{5042}{3}=21{,}350{,}046{,}480$ multisets of three rankings of seven goods; hard profiles with $10\le m\le12$ ($496$), decided by enumeration of $S$ plus a mixed-integer program; the counterexample-guided decision procedure (\path{src/cegar.py}) for $4{,}087$ hard profiles with $13\le m\le21$ (per $m$: $1{,}367$, $950$, $803$, $447$, $242$, $223$, $33$, $22$ profiles for $m=13,14,15,16,17,18,20,21$; largest refuting set $12$); $102$ constructed hard profiles with $m\le36$; refuting sets re-verified without any integer program for $57$ profiles of the searches plus $33$ concatenated ones.

\paragraph{Reproduction.}
The verification directories of the repository each state the result they support. For the finite theorem, the complete verifications are
\begin{verbatim}
cd verification/coverage12
python3 run_verification.py --workers 4
python3 run_tests.py
cd ../coverage15
python3 run_verification.py --m 15 --workers 4
\end{verbatim}
The two \path{run_verification.py} commands compile the C++17 coverage verifiers and check that disjoint partition-index ranges cover all $15{,}400$, respectively $1{,}401{,}400$, partitions; their results are the computer-assisted part of the existence proof. \path{run_tests.py} runs construction and obstruction tests, whose random successes do not extend the theorem's range. The completion and repair tests use exact \texttt{Fraction} arithmetic. Where two implementations exist, they were written independently and agree. The analytic results rely on their proofs; the finite existence theorem additionally relies on exhaustive computation. Neither is formally verified.

%%%%%%%%%%%%%%%%%%%%%%%%%%%%%%%%%%%%%%%%%%%%%%%%%%%%%%%%%%%%%%%%%%%%%%%
\section{Use of Generative AI}\label{app:ai}
This section documents the use of generative AI in the research and in the preparation of the manuscript, as required by the AAMAS 2027 policy. AI systems are not authors. All tools were used through their chat or command-line interfaces.

\begin{center}\small
\begin{tabular}{@{}p{.15\textwidth}p{.2\textwidth}p{.33\textwidth}p{.22\textwidth}@{}}
\toprule
Tool & Version & Role & Prompts\\ \midrule
ChatGPT (OpenAI) & \red{[FILL: model name/version]} & proposed and proved lemmas on the market-aligned conjecture and on covers; first counterexample searches; proof drafts of the completion results & \path{ai_prompts/chatgpt_*.md}\\
Codex (OpenAI) & \red{[FILL: version]} & computation only: extended searches for three-agent counterexamples, checks of the market-aligned conjecture for larger $n,m$ & \path{ai_prompts/codex_*.md}\\
Claude (Anthropic) & Claude Code with claude-sonnet-5-5 for the final audit and rewrite; \red{[FILL: versions of earlier sessions]} & wrote and ran verification code (\path{src/}), the two-agent proof write-up, an independent external review, and the audit and revision of this manuscript & interactive sessions; task-defining prompts not archived verbatim\\
Gemini Deep Research (Google) & \red{[FILL: version]} & initial literature scan that proposed six candidate open problems; every citation and claim was re-checked against the sources, and only Open Question~1 of Barman et al.\ survived & report only; not archived\\
\bottomrule
\end{tabular}
\end{center}

\paragraph{What the systems did.} The systems proposed conjectures, drafted proofs, wrote and ran verification code, and helped edit text. Chat transcripts are not archived; the task-defining prompts given to ChatGPT and Codex are in the directory \path{ai_prompts/} of the supplementary archive (with local file-system paths removed), and the returned proofs and programs were archived in the packages named in the repository's verification directories.

\paragraph{What was verified, and how.}
(1) Every finite claim (counts, enumerations, coverage) was re-run locally from the archived code, and its outputs compared field by field with the archived results. (2) The central finite claims were re-implemented independently of the code that produced them (the two-agent path and monotone checks, the thirteen-good and repair-distance enumerations, the twelve-good disjoint-defect enumeration, the fifteen-good coverage checker, the exact linear-programming cover computations of Appendix~\ref{app:family}). (3) Every proof in the paper was re-derived line by line in a separate audit during the preparation of this submission, and hypotheses were tested on small cases by machine; errors found and corrected in earlier drafts are recorded in the repository's notes. (4) Bibliographic entries were checked against the publishers' or arXiv metadata. No proof assistant was used, and the results are not formally verified.

\end{document}
